% ============================================================
% Abstract (English & Arabic) - Standard UST Spacing Layout
% ============================================================
\thispagestyle{plain}
\vspace*{0.2cm}

% --- English Abstract Header ---
\begin{center}
{\fontsize{15}{18}\selectfont\bfseries ABSTRACT}
\end{center}
\vspace{0.2cm}

\fontsize{11}{16.5}\selectfont
\setstretch{1.45}
This project presents an integrated assistive platform designed for individuals with severe motor disabilities, such as quadriplegia, locked-in syndrome, ALS, and muscular dystrophy, who retain voluntary eye movement. The platform integrates motorized wheelchair navigation and virtual keyboard communication into a unified assistive system controlled via Electrooculography (EOG) signals. The technical architecture comprises a discrete EOG analog front-end with an AD620 instrumentation amplifier, digital filtering in MATLAB (250~Hz sampling rate, moving-average smoothing, and 50~Hz notch filtering), and an Arduino microcontroller with ultrasonic obstacle detection. Preliminary evaluation across five healthy subjects yielded an overall command execution accuracy of 92\% and a mean text entry rate of 16 characters per minute (CPM), demonstrating the operational feasibility of the low-cost prototype.

\vspace{0.8cm}

% --- Arabic Abstract Header ---
\begin{center}
\begin{Arabic}
{\fontsize{15}{18}\selectfont\bfseries المستخلص}
\end{Arabic}
\end{center}
\vspace{0.2cm}

\begin{Arabic}
\fontsize{11.5}{18}\selectfont
\setstretch{1.65}
يُقدم هذا المشروع تصميماً هندسياً لمنصة مساعدة مدمجة تلبي احتياجات الأفراد الذين يعانون من إعاقات حركية شديدة (مثل الشلل الرباعي، متلازمة الانحباس، التصلب الجانبي الضموري ALS، وضمور العضلات) مع احتفاظهم بحركة العين الإرادية. يدمج النظام وظيفتين أساسيتين؛ وهما التحكم في كرسي متحرك بمحركات كهربائية، والتواصل عبر لوحة مفاتيح افتراضية، داخل نظام موحد يعتمد على إشارات حركة العين (Electrooculography - EOG). تعتمد المعمارية التقنية على دائرة تماثلية مخصصة لاكتساب الإشارة وتكبيرها باستخدام مكبر الأجهزة (AD620)، متبوعة بترشيح رقمي عبر بيئة MATLAB بمعدل أخذ عينات 250~Hz ومرشح إيقاف حزمة 50~Hz. أظهرت الاختبارات المعملية الأولية على خمسة متطوعين أصحاء دقة تنفيذ أوامر بلغت 92\% ومتوسط سرعة كتابة 16 حرفاً في الدقيقة، مدعوماً بنظام أمان لتفادي العوائق عبر حساسات الموجات فوق الصوتية، مما يؤكد الجدوى التشغيلية لهذا النموذج الأولي منخفض التكلفة.
\end{Arabic}

\normalsize
\clearpage
